\section{Connection to Classical Superprocess Indices}
\label{app:remarks}

\begin{remark}[The theorem as a superprocess index theorem]\label{rem:whittle}
Adding to an arm a free \emph{keep} transition that pays the running maximum of its collectable
readings would encode the collectable set in the Markovian boxes of \citet{gupta2019markovian},
but it gives the arm two actions and leaves the class of Markov chains for that of decision
processes, which is why the reduction is to the superprocess theorem and not to the Markovian box.
In the language of \citet{whittle1980gittins} and \citet[\S5.2.2]{scully2025gittins}, who call the setting MDP selection, a
nonterminal arm at state $s$ with best previously collectable reading $M$ can \emph{advance}
at reward $-c(s)$ or \emph{keep} $M$ and terminate the selection. When none has been observed,
$M=-\infty$ and keeping is unavailable. Including an outside retirement value $\alpha$ gives
the local value
\[
r+\Phi(s;r)^+,\qquad r=\max\{\alpha,M\},
\]
by Lemma~\ref{lem:nested}. Here advance is compared with the better of \emph{both} immediate
options, not with keeping $M$ alone: retirement can exceed $M$.

Whittle's condition asks for one non-retirement action, chosen independently of $\alpha$,
such that this action or retirement is optimal for every $\alpha$. The choice is
\begin{center}
\begin{tabular}{lll}
\toprule
Condition at the arm state & Retained action & Retirement optimal for \\
\midrule
$M> -\infty$, $M\ge G(s)$ & keep & $\alpha\ge M$ \\
$M<G(s)$ & advance & $\alpha\ge G(s)$ \\
\bottomrule
\end{tabular}
\end{center}
Below the displayed retirement threshold, the retained action is optimal;
at the threshold both may be chosen. Its identity depends only on the fixed
arm state, through $M$ and $G(s)$, while the choice to retire may depend on $\alpha$.
If both $M$ and $G(s)$ are $-\infty$, retirement is optimal for every finite $\alpha$;
any available advance may be the designated action in that vacuous case.
The resulting outside-option threshold is $\max\{M,G(s)\}$. With the incumbent represented
as a deterministic retirement option (equivalently, a dummy arm), it yields
$\pi_{\rm IX}$ in the terminating selection setting of \citet{glazebrook1979stoppable}. This is a connection to classical theory; the direct Bellman
verification in Theorem~\ref{thm:index} already establishes optimality for the stated model.
Nonobligatory-inspection hardness \citep{fu2023nonobligatory} does not apply because an
unobserved reward cannot be kept.
\end{remark}

